% !TEX root = main.tex

\section{\texorpdfstring{Expected Kolmogorov complexity of $\beta$-expansions generated by the Daubechies algorithm}{}}

ENLEVER CETTE SECTION POUR HELMUT, ET FAIRE LA SECTION SUR L'ALGORITHME. POINT.

Bon ALORS: on se rassure.

Let $(\Omega,\AM, \mu)$ be a probability space. 
\begin{definition}
    A $\mu$-test is a measurable function $f : \Omega \to \R$ such that 
    \begin{equation}
        \mu\left\{ \omega \in \Omega : f(\omega) > k \right\} \leq 2^{-k}.
    \end{equation}
\end{definition}
\begin{definition}
    A unit $\mu$-integrable function is a measurable function $f : \Omega \to \R$ such that 
    \begin{equation}
        \int_\Omega f d\mu \leq 1.
    \end{equation}
\end{definition}
\begin{theorem}
    Let $f$ be unit $\mu$-integrable. Then, $\log f$ is a $\mu$-test.
\end{theorem}
\begin{proof}
    Let $f$ be unit $\mu$-integrable. Define, for $k \in \N$,
    \begin{equation}
        A_k := \left\{ \omega \in \Omega :\log f(\omega) > k \right\}.
    \end{equation}
    Suppose that $\log f$ is not a $\mu$-test, i.e., that there exists $k \in \N$ such that $\mu(A_k) > 2^{-k}$. Then,
    \begin{align}
        \int_\Omega f d\mu \geq \int_{A_k} fd\mu > 2^k \mu(A_k) > 1,
    \end{align}
    which is a contradiction.
\end{proof}
\noindent From now on $\Omega = \{0,1\}^\N$. Let $F : x \mapsto x + \log_2(x)$.
\begin{definition}
    Let $A \in \Omega$ and $f_{0|A} : \Omega \to \R$ be defined by 
    \begin{equation}
        f_{0|A}(\omega) := \sup_{n \in \N} \left\{ \log \frac{1}{\mu(\omega_{1:n}\Omega)} - F\circ C[\omega_{1:n}|A] \right\},
    \end{equation}
    for all $\omega \in \Omega$.
\end{definition}
\begin{theorem}
    For all $A \in \Omega$, $f_{0|A}$ is a $\mu$-test.
\end{theorem}
\begin{proof}
    We prove that $2^{f_{0|A}}$ is unit $\mu$-integrable. First, define 
    \begin{equation}
        h_n(\omega) =  \log \frac{1}{\mu(\omega_{1:n}\Omega)} - F \circ C [\omega_{1:n}|A].
    \end{equation}
    Then,
    \begin{align}
        \sum_{n\in \N} \int_\Omega 2^{h_n} d\mu & = \sum_{n \in \N} \int_{\omega \in \Omega} \frac{1}{\mu(\omega_{1:n}\Omega)} 2^{- F \circ C [\omega_{1:n}|A]} d\mu(\omega)\\
        & = \sum_{n \in \N} \sum_{x \in \{0,1\}^n} \frac{1}{\mu(x\Omega)} 2^{- F \circ C [x|A]} \mu(x \Omega)\\
        & = \sum_{n \in \N} \sum_{x \in \{0,1\}^n} 2^{- F \circ C [x|A]} = \sum_{x \in \{0,1\}^\ast} 2^{- F \circ C [x|A]}.
    \end{align}
    For each $x \in \{0,1\}^\ast$, define $x_A \in \{0,1\}^\ast$ to be the canonical string associated to $x$ relatively to $A$, i.e. the lexicographically maximal string satisfying $|x_A| = C(x|A)$. Note that the set $\{ \bar x_A : x \in \{0,1\}^\ast\}$ is prefix-free, and hence by Kraft inequality,
    \begin{equation}
        \sum_{x \in \{0,1\}^\ast} 2^{-\len{\bar x_A}} \leq 1.
    \end{equation}
    Since $\len{\bar x_A} = F \circ C[x|A]$ we get 
    \begin{equation}
        \sum_{n\in \N} \int_\Omega 2^{h_n} = \sum_{x \in \{0,1\}^\ast} 2^{- F \circ C [x|A]} = \sum_{x \in \{0,1\}^\ast} 2^{-\len{\bar x_A}} \leq 1.
    \end{equation}
    We conclude by noting that 
    \begin{align}
        \int_\Omega 2^{f_{0|A}} d\mu = \int_\Omega  2^{\sup_{n\in \N} h_n} d\mu \overset{(a)}{=}  \int_\Omega \sup_{n\in \N} 2^{h_n} d\mu
        \leq \int_\Omega \sum_{n\in \N} 2^{h_n} d\mu
        \overset{(b)}{=} \sum_{n\in \N} \int_\Omega 2^{h_n} d\mu
        \leq 1,
    \end{align}
    where (a) is follows from the fact that $x \mapsto 2^x$ is continuous, and (a) is by the Lebesgue sum-integral inversion theorem.
\end{proof}


In this section, we study the expected Kolmogorov complexity of the $\beta$-expansions generated by Algorithm DAUBECHIES. Recall that on input $s \in I_\beta$, the physical device $T$ satisfies the following input-output relationship
\begin{equation}
    T(s) = \begin{cases}
        0, \ \text{if} \ s \in [0, t - \varepsilon) =: E_T^{(0)},\\
        \text{Randomly} \ 0 \ \text{or} \ 1, \ \text{if} \ s \in [t - \varepsilon, t + \varepsilon] =: S_T,\\
        1, \ \text{if} \ s \in (t + \varepsilon, 1/(\beta-1)] =: E_T^{(1)}.
    \end{cases}
\end{equation}
Note that the sequence $(r_i)_{i \in \{0, \ldots, n\}}$ defined in lines LINES of Algorithm DAUBECHIES is defined by the following recursive relationship
\begin{align}
    r_0 &= s,\\
    r_{i+1} &= \beta r_i - T(r_i) =: K_{\beta,T}(r_i), \ \forall i \in \{0,\ldots, n-1\}.
\end{align}
Therefore, the output of Algorithm DAUBECHIES $\mathcal A_\beta(s,n,T)$ on input $s \in [0,1]$ and $n \in \N$, is expressed as
\begin{equation}
    \mathcal A_\beta(s,n,T) = \coprod_{i=0}^{n-1} T\left(K_{\beta,T}^i(s)\right).
\end{equation}
In order to formalize better the problem, we can define $T(s)$ formally as a random variable on $\{0,1\}$, i.e., as a function $T(s) : \{0,1\} \to \{0,1\}$ satisfying
\begin{equation}
    T(s)(x) := \begin{cases}
        0, \ \text{if} \ s \in  E_T^{(0)},\\
        x, \ \text{if} \ s \in  S_T,\\
        1, \ \text{if} \ s \in  E_T^{(1)},
    \end{cases}
    \ \forall x \in \{0,1\}.
\end{equation}
Since the generation of the sequence $(r_i)_{i \in \N}$ might rely on an arbitrarily large number of Bernoulli experiments, we define an extension $\tilde T(s)$ of $T(s)$, that will be useful to consider Algorithm DAUBECHIES as a random variable. Namely, we define $\tilde T(s) : \{0,1\}^\N \to \{0,1\}$ by 
\begin{equation}
    \tilde T(s)(\xf) := \begin{cases}
        0, \ \text{if} \ s \in  E_T^{(0)},\\
        \xf_1, \ \text{if} \ s \in  S_T,\\
        1, \ \text{if} \ s \in  E_T^{(1)},
    \end{cases}
    \ \forall \xf \in \{0,1\}^\N.
\end{equation}
Thanks to the formulation above, we can see Algorithm DAUBECHIES as first generating an infinite sequence $\xf \in \{0,1\}^\N$ out of infinitely many Bernoulli experiments, and then reading successively the outcomes of this experiment, i.e., at each time step, the algorithm reads $\xf_1$, and then discards it, such that the sequence $\xf$ is transformed into the sequence $\yf := \xf_2\xf_3\ldots$. This is usually denoted by $\yf = \sigma(\xf)$, where $\sigma : \{0,1\}^\N \to \{0,1\}^\N$ is called the left-shift. Then, in the next time step, the algorithm will read $\yf_1 = \xf_2$ to generate the next bit of the $\beta$-expansion of $s$.  Note that when $s \notin S$, then the generated bit does not depend on the Bernoulli experiments, as it is generated deterministically. Hence, the Algorithm DAUBECHIES does not need to discard $\xf_1$, and wait until the next time step for which $s \in S$ to discard it. Accordingly, at each time step, the function $\sigma_s : \{0,1\}^\N \to \{0,1\}^\N$ defined by 
\begin{equation}
    \sigma_s(\xf) := \begin{cases}
        \xf \ \text{if} \ s \notin S,\\
        \sigma(\xf), \ \text{if} \ s \in S,
    \end{cases} \ \forall s \in I_\beta, \ \xf \in \{0,1\}^\N,
\end{equation}
is applied to the infinite sequence $\xf$ of Bernoulli experiments. We are now ready to state formally the definition of the Algorithm DAUBECHIS as a random variable. First, we model $K_{\beta,T}(s)$ as a random variable on $\{0,1\}^\N$, i.e., as a function $K_{\beta,T}(s) : \{0,1\}^\N \to I_\beta \times \{0,1\}^\N$.  defined by 
\begin{equation}
    K_{\beta,T}(s)(\xf) := (\beta s - \tilde T(s)(\xf_1), \sigma_s(\xf)),
    \ \forall \xf \in \{0,1\}^\N.
\end{equation}
In order to simplify the notations, we might consider $\tilde T$ as a function from $I_\beta \times \{0,1\}^\N$ into $\{0,1\}$ and $K_{\beta,T}$ as a function from $I_\beta \times \{0,1\}^\N$ into itself, by making the confusion between $T(s)(\xf)$ and $T(s,\xf)$, and the confusion between $K_{\beta,T}(s)(\xf)$ and $K_{\beta,T}(s,\xf)$. This simpler formalism allows to define the output of the Algorithm DAUBECHIES as a random variable as well, i.e., as a function $\mathcal A_\beta(s,n,T) : \{0,1\}^\N \to \{0,1\}^n$, by 
\begin{equation}
    \mathcal A_\beta(s,n,T)(\xf) = \coprod_{i=0}^{n-1} \tilde T\left(K_{\beta,T}^i(s,\xf)\right), \ \forall \xf \in \{0,1\}^\N.
\end{equation}
We repeat that Algorithm DAUBECHIES can generate different outputs of the same input $s \in [0,1]$. Namely, for $s \in [0,1]$ and $n \in \N$, we denote by $\Sigma_\beta(s,n,T) \subseteq \{0,1\}^n$ the set of possible outputs of Algorithm DAUBECHIES on input $(s,n)$. Note that 
\begin{equation}
    \Sigma_\beta(s,n,T) = \left\{ \mathcal A_\beta(s,n,T)(\xf) : \xf \in \{0,1\}^\N\right\} = \mathcal A_\beta(s,n,T)\left( \{0,1\}^\N\right).
\end{equation} 
We let also $\NN_\beta(s,n,T) :=  \# \left(\Sigma_\beta(s,n,T) \right)$. For the specific case of $T$ satisfying $E_T^{(0)} = [0, \beta^{-1})$ and $S_T = [\beta^{-1}, (\beta(\beta-1))^{-1}]$, $\NN_\beta(s,n,T)$ is a well-studied quantity usually called the \textit{growth rate for $\beta$-expansions}, as it counts how many $\beta$-expansions the number $s$ has. We now leverage the formalism introduced above to establish lower bounds on the expected Kolmogorov complexity of the $\beta$-expansions it generates. Recall that for a fixed number $\beta \in (1,2)$, we have defined the equivalence relationship $\sim_\beta$ over $\{0,1\}^\ast$ by 
\begin{equation}
    x \sim_\beta y \iff n := |x| = |y| \ \text{and} \ \sum_{i=1}^n x_i \beta^{-i} = \sum_{i=1}^n y_i \beta^{-i}, \ \forall x,y \in \{0,1\}^\ast.
\end{equation}
Moreover, we have defined the equivalence class $[x]_\beta := \{y \in \{0,1\}^\ast : x \sim_\beta y\}$, for all $x \in \{0,1\}^\ast$, which, seen as a multivalued function $[\cdot]_\beta : \{0,1\}^\ast \rightrightarrows \{0,1\}^\ast$, is proven to be computable in the proof of Corollary \ref{cor:M is computable for algebraic beta}.We define 
\begin{equation}
    \tilde\NN_\beta(s,n,T) := \# \left(\Sigma_\beta(s,n,T) / \sim_\beta\right).
\end{equation}
Note that if $\beta$ is not algebraic, then $[x]_\beta = \{x\}$ for every $x \in \{0,1\}^n$, hence $\tilde\NN_\beta(s,n,T) = \NN_\beta(s,n,T)$. We now construct a multivalued function $f_{\beta, 1 \to all} : \{0,1\}^\ast \rightrightarrows \{0,1\}^\ast$ that are computable relatively to $\beta$, such that
\begin{equation}
    \langle \Sigma_\beta(s,n,T) \rangle \in f_{\beta, 1 \to all}\left(x\right). 
\end{equation}
To build $f_{\beta, 1 \to all}$, we first follow the same intuitions that lead to the construction of $f_{\beta \to 2}$ and $f_{2 \to \beta}$ in the PREVIOUS SECTIONS. Namely, the knowledge of $x \in \Sigma_\beta(s,n,T)$ allows to infer that
\begin{equation}
    s \in I(x) := \left[\sum_{i=1}^{n} x_i \beta^{-i},  \sum_{i=1}^{n} x_i \beta^{-i}+ \frac{\beta^{-n}}{\beta-1} \right],
\end{equation}
and further deduce that 
\begin{equation}
    \sum_{i=1}^n y_i \beta^{-i} \in J(x) := \left[\sum_{i=1}^{n} x_i \beta^{-i} - \frac{\beta^{-n}}{\beta-1},  \sum_{i=1}^{n} x_i \beta^{-i}+ \frac{\beta^{-n}}{\beta-1} \right],
\end{equation}
for all $y \in \Sigma_\beta(s,n,T)$. Hence, by defining $g_\beta : \{0,1\}^\ast \rightrightarrows \{0,1\}^\ast$ as 
\begin{equation}
    g_\beta(x) := \left\{ y \in \{0,1\}^n : n = |x|,  \sum_{i=1}^n y_i \beta^{-i} \in J(x)\right\},
\end{equation}
for all $x \in \{0,1\}^\ast$, we get 
\begin{equation}
    \Sigma_\beta(s,n,T) \subseteq g_\beta(x), \ \forall x \in \Sigma_\beta(s,n,T).
\end{equation}
Now, for $x \in \{0,1\}^\ast$, define $\tilde g_\beta(x) := g_\beta(x)/ \sim_\beta$, and $\tilde \NN_g(x) := \# \tilde g_\beta(x)$. Note that the function defined by
\begin{equation}
    \begin{array}{rccc}
        \varphi_\beta : & \{0,1\}^n / \sim_\beta & \to & I_\beta\\
        & [x]_\beta & \mapsto & \sum_{i=1}^n x_i \beta^{-i}
    \end{array}
\end{equation}
is injective into $I_\beta$. We define an ordering $\prec$ on $\{0,1\}^n / \sim_\beta$ by 
\begin{equation}
    u \prec v \iff \varphi_\beta(u) < \varphi_\beta(v), \ \forall u,v \in \{0,1\}^n / \sim_\beta.
\end{equation}
For $x \in \{0,1\}^\ast$, we enumerate the elements of $\tilde g_\beta(x)$ as follows. We let $u_1(x), \ldots, u_{k}(x)$ such that $\{u_1(x), \ldots, u_k(x)\} = \tilde g_\beta(x)$ and $i < j \Rightarrow u_i(x) \prec u_j(x)$, for all $i,j \in \{1, \ldots, k\}$. We can then express $\Sigma_\beta(s,n,T)$ simply in function of the elements of $\tilde g_\beta(x)$. 
\begin{lemma}
    Let $\beta \in (1,2)$, $s \in [0,1]$, $n \in \N$, $T$ be a physical device and $x \in \Sigma_\beta(s,n,T)$. Then, there exists $i,j \in \{1, \ldots, \tilde \NN_g(x)\}$, $i \leq j$, such that 
    \begin{equation}
        \Sigma_\beta(s,n,T) = \bigcup_{\ell=i}^ju_\ell(x).
    \end{equation}
\end{lemma}
\begin{proof}
    Let $\beta \in (1,2)$, $s \in [0,1]$, $n \in \N$, $T$ be a physical device and $x \in \Sigma_\beta(s,n,T)$. The proof is in two parts: we first prove that EH NON IL EST PAS VRAI CE THEOREME BOLOSS.
\end{proof}


For a fixed $s \in [0,1]$, $n \in \N$ and $\xf \in \{0,1\}^\N$, recall that the sequence $(r_i)_{i \in \N}$ results from iterations of the function $K_{\beta,T}$, and define
\begin{equation}
    h_\beta(s,n,\xf,T) := \# \{r_i \in S : i \in \{1,\ldots,n  \} \}.
\end{equation}
In fact, $h_\beta(s,n,\xf,T)$ counts the number of times that Algorithm DAUBECHIES had to read the first bit of the infinite sequence of Bernoulli experiments in order to generate its first bit, i.e., the Algorithm had to read exactly the $h_\beta(s,n,\xf,T)$ first bits of the infinite sequence $\xf \in \{0,1\}^\N$ in order to generate its output. Note that the output of the Algorithm does not depend on the bits of $\xf$ located after the position $h_\beta(s,n,\xf,T)$, hence, for every $x \in \Sigma_\beta(s,n,T)$, there exists a sequence $\xf \in \{0,1\}^\N$ such that 
\begin{equation}
    \mathcal A_\beta(s,n,T)\left(\xf_{1:h_\beta(s,n,\xf,T)} \{0,1\}^\N\right) = \{x\}.
\end{equation}
Accordingly, we define 
\begin{equation}
    \noise_\beta(s,n,T,x) := \left\{u \in \{0,1\}^{\leq n} : \mathcal A_\beta(s,n,T)\left(u \{0,1\}^\N\right) = \{x\}\right\},
\end{equation}
for all $x \in \Sigma_\beta(s,n,T)$. We first construct two multivalued functions $f_{\beta, 1 \to all}, f_{\beta, 1+all \to \noise} : \{0,1\}^\ast \rightrightarrows \{0,1\}^\ast$ that are computable relatively to $\beta$, such that
\begin{equation}
    \langle \Sigma_\beta(s,n,T) \rangle \in f_{\beta, 1 \to all}\left(x\right), 
\end{equation}
and
\begin{equation}
    \langle \noise_\beta(s,n,T,x) \rangle \in f_{\beta , 1+all \to \noise}\left(\langle x,\langle \Sigma_\beta(s,n,T) \rangle \rangle \right),
\end{equation}
for all $s \in [0,1]$, $n \in \N$, and $x \in \Sigma_\beta(s,n,T)$. 



This implicitely defines a multivalued function $f_{\beta, 1 \to \noise} : \{0,1\}^\ast \rightrightarrows \{0,1\}^\ast$ by 
\begin{equation}
    f_{\beta, 1 \to \noise}(x) := f_{\beta , 1+all \to \noise}\left(\langle x, f_{\beta, 1 \to all}\left(x\right) \rangle \right), \ \forall x \in \{0,1\}^\ast.
\end{equation}




We denote by  the se
Note also that, in the above equation, the sequence $\xf_{1:h_\beta(s,n,\xf,T)}$ is uniquely determined by $x$.



Hence, the generation of the $\beta$-expansion can be hence seen as a nondeterministic process, where the nondeterministic part occurs only when the variable $r$ in Algorithm 4 belongs to the interval $[t - \varepsilon, t + \varepsilon]$. Indeed, when $r \in [t - \varepsilon, t + \varepsilon]$, $T$ delivers a random value, that can be seen as introducing noise. Given a real number $s \in [0,1]$, we let Algorithm 4 run for infinitely many time step, i.e. $n = \infty$. Denote by $i_1, i_2, \ldots $ the time steps for which the variable $r \in [t - \varepsilon, t + \varepsilon]$, and denote by $\xf \in \{0,1\}^\N$ be the sequence generated by the Algorithm, which is a $\beta$-expansion of $s$. The sequence $\xf$ is then generated deterministically, except for those elements $\xf_{i_1}, \xf_{i_2}, \ldots$ that have been generated randomly. We will show that, indeed, the sequence $\xf_{i_1}\xf_{i_2}\ldots $ can be seen as a noise, because the difference between the algorithmic complexity of $\xf$ and the algorithmic complexity of the $\beta$-expansion $\yf$ of $s$ that would be generated with a perfect threshold $t$ is lower bounded by the algorithmic complexity of $\xf_{i_1}\xf_{i_2}\ldots$. This is formally expressed as 
\begin{equation}
    \underline{\Delta}_\beta(\xf) \geq \log_\beta(2) \liminf_{n \to \infty} \frac{K[\xf_{i_1} \ldots \xf_{i_n}]}{n}\cdot
\end{equation}
OKKK: ALORS C'EST LE GROS BORDEL, LA FONCTION F 2 to BETA ETAIT PAS CALCULABLE !!
Accordingly, we define a function that converts a nondeterministic $\beta$-expansion of $s$ into the corresponding deterministic $\beta$-expansion and the noise introduced.
\begin{definition}\label{def:multivalued function converting from nondeterministic beta to deterministic beta + noise}
    OK JE POURRAIS D'ABORD DESIGNER LA FONCTION SEULEMENT POUR LE NOISE, PUIS LA COMBINER AVEC CELLE POUR L'EXPANSION BINAIRE
    Bon, cette fonction, elle prend comme input 2 entiers en plut de $x$. Ces entiers donnent la position exacte du premier et du dernier cluster de $\beta$-expansions de $s$. Ensuite, on peut donc extraire toutes les $\beta$-expansions de $s$. Là, on peut dessiner un arbre des beta-expansions et placer notre beta expansion dans cet arbre, et extraire chaque bit qui correspond à un branching de l'arbre. Formellement,
    Let $f^\beta_{ND\to D} : \{0,1\}^\ast \rightrightarrows \{0,1\}^\ast$ be the multivalued function defined as 
\begin{align}\label{eq:main:def:multivalued function converting from base beta to base 2 + noise}
    f^\beta_{ND\to D}(\langle x, k, \ell \rangle) &:= x_{i_1, \ldots, i_h},
\end{align}
for all $x \in \{0,1\}^{n\log_\beta(2)}$ for some $n \in \N$, and $f^\beta_{ND\to D}(x) = \epsilon$ for $x \in \{0,1\}^m$ where $m \neq n \log_\beta(2)$ for all $n \in \N$.
\end{definition} 
PROUVER LE RESULTAT POUR $s \in \Q$??




\subsection{Formalisation of the ergodic problem}
Suppose that on input $s \in I_\beta$, the physical device $T$ delivers a bit $0$ or $1$, according to a Bernoulli probability distribution parametrized by the input $s$. We denote this parameter by $p(s)$. A example of $p(s)$ would be the following figure. For $n \in \N$ and $s \in I_\beta$, we denote by $\mathcal A_p(s,n) \in \{0,1\}^n$ the random variable that models the output of Algorithm 4 on input $s \in I_\beta$ and $n \in \N$, depending on $p$. Moreover, we define $\mathcal A_p(s) \in \{0,1\}^\N$ the be the random variable that models the output of Algorithm 4, it we would let it run infinitely long. More formally, this is expressed as follows. Then,
\begin{equation}
    \begin{array}{rccc}
        T_p : & I_\beta \times [0,1]^\N & \to & \{0,1\}\\
                       & r, \tff & \mapsto & \displaystyle  
                       \begin{cases}
                        0 \text{ if } \tff_1 \leq p(r)\\
                        1 \text{ if } \tff_1 > p(r).
                       \end{cases},
    \end{array}
\end{equation}
and the following equation of evolution
\begin{equation}
    \begin{array}{rccc}
        \mathcal K_p : & I_\beta \times [0,1]^\N & \to & I_\beta \times [0,1]^\N\\
                       & r, \tff &\displaystyle \mapsto &  \left(\beta r - T_p(r,\tff_1), \sigma \tff\right)
    \end{array}
\end{equation}
This translates to 
\begin{equation}
    \mathcal A_p(s, \tff) := \coprod_{i=1}^\infty T_p\left( \mathcal K_p^{i-1}(s, \tff)\right) =: \coprod_{i=1}^\infty a_i(s, \tff).
\end{equation}
Define $\supp p := \overline{\{s \in I_\beta : p(s) \in (0,1)\}}$, and for $s, \tff$, let $(I_i(s,\tff))_{i \in \N}$ be the increasing sequence of natural numbers satisfying 
\begin{equation}
    \{I_i(s,\tff) : i \in \N\} = \{ j \in \N : a_j(s,\tff) \in \supp p\}.
\end{equation}
Then, we can define the sequence of elements that where generated by a random process as 
\begin{equation}
    \mathcal R_p(s, \tff) := \coprod_{i=1}^\infty a_{I_i(s,\tff)}(s, \tff) := \coprod_{i=1}^\infty r_i(s, \tff).
\end{equation}

% Given a fixed number $s \in I_\beta$, we write $\mathcal A_s(\tff) := \mathcal A(s,\tff)$. Then, for $s \in I_\beta$, 
$A(s,\tff)$ is a random variable of distribution
\begin{equation}
    \mathcal A_{\#} \left(\lambda_\beta \otimes\lambda^\N\right),
\end{equation}
where $\lambda^\N$ is the infinite product of the Lebesgue measure, such that 
\begin{equation}
    \lambda^\N(C_1\times C_2 \times \ldots \times C_n \times [0,1]^\N) = \prod_{i=1}^n \lambda(C_i),
\end{equation}
for all $C_1, \ldots, C_n \subseteq [0,1]$ Borel sets. Then, for all $n \in \N$, 
\begin{equation}
    \mathbb E[K[\mathcal A_p(s, n)]] = \int_{\{0,1\}^\N}  K[\xf_{1:n}] d\mathcal A_{\#}\left(\lambda_\beta \otimes\lambda^\N\right).
\end{equation}



\subsection{Ergodic properties of the Daubechies algorithm: piecewise-constant case}

\subsection{Study for the piecewiese-constant case}

\subsection{A conjecture for the general case}

